\documentclass{beamer}


\usepackage[utf8]{inputenc}
\usepackage[spanish]{babel}
\usepackage[T1]{fontenc}
\usepackage{graphicx}
\usepackage{tikz}
\usepackage{xcolor}

\title{Python}
\author{
Guillem Borrell i Nogueras\\
\texttt{guillemborrell@gmail.com}
}

\begin{document}

\begin{frame}
\begin{center}
 \includegraphics[width=9cm]{files/python-logo-generic.pdf}\\
 % python-logo-generic.pdf: 389x115 pixel, 72dpi, 13.72x4.06 cm, bb=0 0 389 115
Guillem Borrell i Nogueras\\

ETSIA, Octubre 2007
\end{center}

\end{frame}


\begin{frame}
 \frametitle{Preguntas...}
 \begin{itemize}
 \item ¿Por qué se llama Python?
 \item ¿Quién usa Python?
 \item ¿Para qué sirve Python?
 \begin{itemize}
  \item Principales características de Python
 \end{itemize}
 \item ¿Por qué Python se está volviendo tan popular?
 \begin{itemize}
  \item \emph{Zen of Python}
 \end{itemize}
 \item ¿Por qué Python y no otro lenguaje?
 \begin{itemize}
  \item Ejemplos
 \end{itemize}

 \item ¿Qué puede ofrecer Python al HPC?
 \end{itemize}
\end{frame}

\begin{frame}
 \frametitle{¿Por qué se llama Python?}
 Python debe su nombre a...
\end{frame}

\begin{frame}
 \begin{center}
 \includegraphics[width=9cm]{files/python.jpg}\\
 % python.jpg: 760x544 pixel, 72dpi, 26.81x19.19 cm, bb=0 0 760 544
Monty Python
\end{center}

\end{frame}

\begin{frame}
 \frametitle{¿Quién usa Python?}

\begin{center}
 \includegraphics[width=3cm]{files/google.jpg}
 \includegraphics[width=3cm]{files/Industrial_Light_and_Magic.jpg}
 \includegraphics[width=3cm]{files/NASA_Logo.jpg}
 \includegraphics[width=3cm]{files/honeywell_logo.jpg}
\end{center}
\end{frame}

\begin{frame}
 \frametitle{¿Para qué sirve Python?}
 \begin{center}
 Prácticamente para cualquier cosa que se nos pueda ocurrir
 \end{center}

\end{frame}

\begin{frame}
 \frametitle{Desde páginas web...}
\begin{center}
 \includegraphics[width=9cm]{files/snapshot1.jpg}\\
 % snapshot1.jpg: 925x659 pixel, 762dpi, 3.08x2.20 cm, bb=0 0 87 62
Plone, Zope
\end{center}

\end{frame}

\begin{frame}
 \frametitle{...A utilidades para bioinformática}
\begin{center}
 \includegraphics[width=8cm]{files/pymol.jpg}\\
 % pymol.jpg: 720x540 pixel, 72dpi, 25.40x19.05 cm, bb=0 0 720 540
Pymol
\end{center}

\end{frame}

\begin{frame}
 \frametitle{Principales características de Python}

\begin{itemize}
 \item Software Libre (Licencia estilo BSD)
 \item Interpretado
 \item Interactivo
 \item Multiparadigma
 \begin{itemize}
  \item Procedimental
  \item Modular
  \item Orientado a Objetos
 \end{itemize}
 \item Multiplataforma
 \item Especificación especialmente corta (IronPython, Jython, PyPy, Stackless)
 \item \emph{Incluye las pilas}
 \item ...
\end{itemize}

\end{frame}

\begin{frame}
 \frametitle{¿Por qué Python se está volviendo tan popular}
\begin{itemize}
 \item Fácil de aprender
 \item Fácil de ampliar
 \item Consistente por diseño
 \item Impone un buen estilo de programación
 \item Soporta todas las prácticas propuestas por XP, Agile.
\end{itemize}
\begin{center}
 Porque es divertido
\end{center}

\end{frame}

\begin{frame}[containsverbatim]
\frametitle{\emph{Zen of Python, by Tim Peters}}

\begin{verbatim}
>>> import this
\end{verbatim}

\begin{itemize}
 \item Beautiful is better than ugly.
 \item Explicit is better than implicit.
 \item Simple is better than complex.
 \item Complex is better than complicated.
 \item Flat is better than nested.
 \item Sparse is better than dense.
 \item Readability counts.
 \item Special cases aren't special enough to break the rules.
 \item Although practicality beats purity.
\end{itemize}
\end{frame}

\begin{frame}
\begin{itemize}
 \item Errors should never pass silently.
 \item Unless explicitly silenced.
 \item In the face of ambiguity, refuse the temptation to guess.
 \item \textcolor{red}{There should be one --and preferably only one-- obvious way to do it.}
 \item Although that way may not be obvious at first unless you're Dutch.
 \item Now is better than never.
 \item Although never is often better than *right* now.
 \item If the implementation is hard to explain, it's a bad idea.
 \item If the implementation is easy to explain, it may be a good idea.
 \item Namespaces are one honking great idea -- let's do more of those!

\end{itemize}



\end{frame}


\end{document}
